\section{Plain-Language Protocol Summary}

\subsection{The study in one paragraph}

This is a first-in-human study of two new things used together: a drug called
daraxonrasib, given before and after surgery, and an eight-arm robotic system
that performs the surgery under the continuous control of a human surgeon, with
an on-premises computer model advising that surgeon. The operation is a
pancreaticoduodenectomy, usually called a Whipple, which removes the head of the
pancreas along with the first part of the small intestine, part of the bile
duct, and the gallbladder, and then rebuilds three connections. The disease is
pancreatic ductal adenocarcinoma, PDAC, carrying a KRAS G12 mutation. PDAC is
the third leading cause of cancer death in the United States, and five-year
survival across all stages is about 13 percent \cite{Siegel2025}. The primary
question this study asks is not whether the combination cures the disease. It is
whether the combination is safe enough, and practical enough, to be tested in a
study that could answer that. Up to 18 participants will take part.

\subsection{The words you will hear, defined once}

The consent conversation moves faster if the vocabulary is already settled. The
terms below are the ones this protocol cannot avoid using. Each is defined here
in the sense the protocol uses it, which is occasionally narrower than the sense
the same word carries in general oncology.

\vspace{0.30em}

\begin{xltabular}{\textwidth}{L{3.9cm} Y}
\toprule
\textbf{Term} & \textbf{What it means in this protocol} \\
\midrule
\endfirsthead
\multicolumn{2}{@{}l}{\footnotesize\itshape Table continued from the previous page.}\\
\toprule
\textbf{Term} & \textbf{What it means in this protocol} \\
\midrule
\endhead
\midrule
\multicolumn{2}{r@{}}{\footnotesize\itshape Table continued on the next page.}\\
\endfoot
\bottomrule
\endlastfoot
PDAC & Pancreatic ductal adenocarcinoma, the most common form of pancreatic cancer and the only one this study enrolls \\
Whipple & Pancreaticoduodenectomy: removal of the head of the pancreas, the duodenum, the gallbladder and bile duct, and part of the stomach, followed by three reconnections \\
Resectable & The tumour can be removed with a clear margin without reconstructing a major artery. Borderline and locally advanced disease are not eligible \\
KRAS G12 & A family of mutations in the KRAS gene. Daraxonrasib targets this family, so the mutation must be confirmed on tissue before enrollment \\
Daraxonrasib (RMC-6236) & The investigational oral drug in this study, a multi-selective inhibitor of the RAS pathway \\
Physical AI & Software that proposes actions for a machine that touches the patient. Here it proposes; it never actuates \\
On-premises & The language model runs on hardware inside the hospital. No patient data leaves the building during a procedure \\
Advisory & A proposed motion or step produced by the model. An advisory is not an instruction until a surgeon signs it \\
Safety gate & Deterministic software that tests every advisory against the numeric limits before a surgeon ever sees it \\
No-fly corridor & A 4 mm exclusion volume around the superior mesenteric and portal veins that no instrument tip may enter \\
Heartbeat bus & A 10 kHz signal every arm must hear. Silence for one interval triggers an automatic halt \\
Emergency stop & The surgeon-initiated halt: at most 3 ms across arms, at most 500 ms system-wide \\
Phase 0 gate & The simulation-validation milestone that must be signed before any patient-facing robotic activity \\
USL & Unified Safety Level, the composite score the Phase 0 gate must reach, at least 7.0 \\
3+3 escalation & The dose-finding rule: three participants per level, expandable to six, with a level opening only when the level below is cleared \\
Sentinel window & The observation period after the first participant at a dose level, which must close before the second is dosed \\
DLT & Dose-limiting toxicity, a prespecified severe drug reaction that counts against the 3+3 rule \\
ISGPS grade & The international grading of a pancreatic fistula, the most common serious complication of a Whipple \\
R0 & A resection margin with no tumour at the inked edge, assessed on the day-90 pathology report \\
PFS and OS & Progression-free survival and overall survival, followed here to 24 months and reported as descriptive, not comparative \\
\end{xltabular}

\vspace{0.30em}

\subsection{What is being tested, in three parts}

\textbf{The drug.} Daraxonrasib, also called RMC-6236, is a RAS(ON)
multi-selective inhibitor: it binds the active form of the RAS protein that a
KRAS G12 mutation leaves switched on. It has been studied in previously treated
metastatic pancreatic cancer in a randomised Phase 3 programme
\cite{rasolute302, oreilly2026darax}. In this study it is used differently, in
the perioperative setting, with a planned pause before the operation and an
advised restart at 7, 14, or 21 days afterwards, keyed to how the pancreatic
join is healing and to the drug level in the blood.

\textbf{The operation.} A robotic Whipple is an established operation at
high-volume centres. A nationwide series of 1000 robotic
pancreatoduodenectomies shows that the rate of a serious pancreatic leak and the
rate of death within 90 days both fall substantially as an institution's
experience grows \cite{DutchCohort2025}. What is new here is not the robot: it
is the eight-arm configuration, the hard force and geometry limits enforced in
hardware, and the advisory layer above them \cite{onpremwhippl}.

\textbf{The software.} An on-premises large language model observes the sensor
stream and proposes the next operative step. It cannot move anything. A
systematic review of autonomy levels in FDA-cleared surgical robots found that
no cleared system exceeds task-level autonomy under supervision
\cite{LeeAuto2024X}, and this system does not either. Section 7.3 sets out
exactly where the model sits in the chain, and Figure 19 draws it message by
message.

\subsection{Your journey, start to finish}

Figure 3 draws the whole journey in five phases: screening of up to 28 days, a
baseline day, the operation on day 0, an acute window of days 1 to 7, and
follow-up through day 30, day 90, and every 12 weeks to the 24-month overall
survival endpoint. The figure is drawn so that every node the participant
controls is filled Corporate Blue and every node the sponsor controls is left
white. Four of the fifteen nodes are the participant's. The figure does not
pretend that is a majority; it makes the balance visible so that it can be
argued about rather than assumed.

The structure of the journey is adapted from the author's simulation of a fully
autonomous, single-patient journey through a regulated Physical AI oncology
trial \cite{paipatientjr}. That simulation was non-small cell lung cancer, and
four differences matter to a reader with PDAC. The operation is a
pancreaticoduodenectomy with three anastomoses rather than a lobectomy with one
bronchial closure. The drug is a RAS(ON) inhibitor rather than a checkpoint
inhibitor. The dominant early complication is a postoperative pancreatic fistula
rather than a prolonged air leak. And the survival baseline is far lower, which
changes what a participant is weighing when they decide.

\begin{pafloat}
\begin{pafig}[mermaid-type: flowchart LR, full page]
\node[mmin,text width=30mm]  (a1) at (0,0)      {Referral and eligibility review};
\node[mmin,text width=30mm]  (a2) at (0,-2.2)   {KRAS G12 confirmation on tissue};
\node[mmin,text width=30mm]  (a3) at (0,-4.4)   {Staging: resectable or borderline};
\node[mmgoal,text width=32mm](a4) at (0,-6.8)   {\textbf{You decide}\\consent, or decline with no consequence};
\node[mmin,text width=30mm]  (b1) at (5.4,0)    {Phase 0 sign-off, $\geq 1000$ runs, $\geq 2$ frameworks};
\node[mmin,text width=30mm]  (b2) at (5.4,-2.4) {USL $\geq 7.0$ and safety matrix verified};
\node[mmgoal,text width=32mm](b3) at (5.4,-5.0) {\textbf{You decide}\\Physical AI opt-out still open here};
\node[mmmid,text width=30mm] (c1) at (10.8,0)   {Eight-arm robotic Whipple, continuous oversight};
\node[mmin,text width=30mm]  (c2) at (10.8,-2.4){Force caps $\leq 3$ N per arm, $\leq 18$ N summed};
\node[mmin,text width=30mm]  (c3) at (10.8,-4.8){Vascular no-fly gating, SMV and portal vein};
\node[mmin,text width=30mm]  (d1) at (16.2,0)   {ISGPS fistula grading};
\node[mmin,text width=30mm]  (d2) at (16.2,-2.2){Daraxonrasib restart advisory, T+7 / T+14 / T+21};
\node[mmgoal,text width=32mm](d3) at (16.2,-4.8){\textbf{You decide}\\accept or decline the advised restart};
\node[mmin,text width=30mm]  (e1) at (21.6,0)   {Day 30, Clavien-Dindo, primary safety window};
\node[mmin,text width=30mm]  (e2) at (21.6,-2.4){Day 90, R0 pathology and 90-day mortality};
\node[mmin,text width=30mm]  (e3) at (21.6,-4.8){Every 12 weeks to the 24-month endpoint};
\node[mmgoal,text width=32mm](e4) at (21.6,-7.2){\textbf{You decide}\\withdraw at any visit, care continues};
\begin{scope}[on background layer]
\node[mmlane,fit=(a1)(a2)(a3)(a4)] (L1) {};
\node[mmlane,fit=(b1)(b2)(b3)]     (L2) {};
\node[mmlane,fit=(c1)(c2)(c3)]     (L3) {};
\node[mmlane,fit=(d1)(d2)(d3)]     (L4) {};
\node[mmlane,fit=(e1)(e2)(e3)(e4)] (L5) {};
\end{scope}
\node[mmlanetitle,anchor=south west] at (L1.north west) {Screening, up to 28 days};
\node[mmlanetitle,anchor=south west] at (L2.north west) {Baseline, day $-1$};
\node[mmlanetitle,anchor=south west] at (L3.north west) {Surgery, day 0};
\node[mmlanetitle,anchor=south west] at (L4.north west) {Acute, days 1 to 7};
\node[mmlanetitle,anchor=south west] at (L5.north west) {Follow-up, to 24 months};
\draw[mmedge] (a1) -- (a2); \draw[mmedge] (a2) -- (a3); \draw[mmedge] (a3) -- (a4);
\draw[mmedge] (b1) -- (b2); \draw[mmedge] (b2) -- (b3);
\draw[mmedge] (c1) -- (c2); \draw[mmedge] (c2) -- (c3);
\draw[mmedge] (d1) -- (d2); \draw[mmedge] (d2) -- (d3);
\draw[mmedge] (e1) -- (e2); \draw[mmedge] (e2) -- (e3); \draw[mmedge] (e3) -- (e4);
\draw[mmedgeb] (L1.east) -- (L2.west);
\draw[mmedgeb] (L2.east) -- (L3.west);
\draw[mmedgeb] (L3.east) -- (L4.west);
\draw[mmedgeb] (L4.east) -- (L5.west);
\node[mmgray,text width=46mm] (out) at (10.8,-11.8) {Standard-of-care pathway, unchanged and available at every exit};
\draw[mmedged] (a4.south) -- (0,-10.2) -- node[mmlabel,above,pos=0.42]{declining ends it here} (out.west |- 0,-10.2) -- (out.north west);
\draw[mmedged] (b3.south) -- (5.4,-10.2) -- (out.north);
\draw[mmedged] (e4.south) -- (21.6,-10.2) -- node[mmlabel,above,pos=0.42]{withdrawal at any time} (out.east |- 0,-10.2) -- (out.north east);
\node[font=\scriptsize\sffamily\itshape,text=protogray,anchor=north west,align=left]
  at (-2.4,-13.0) {Corporate Blue marks a node the participant controls; white marks a node the sponsor controls; the single lighter-blue node\\is the one investigational act. Four of the fifteen nodes are the participant's, and the figure says so rather than implying more.};
\end{pafig}
\vspace{-0.7cm}
\figcaption{Figure 3. The participant journey across five phases, screening through the 24-month\\
endpoint, with every node the participant controls filled Corporate Blue and every\\
node the sponsor controls left white, so that the balance of control is visible.}
\end{pafloat}

\subsection{Your schedule, and what it costs you in time}

The parent protocol carries a Schedule of Activities: a matrix of activities
against visits, with a cross in each cell where something is collected. It is
the right document for a study coordinator and the wrong one for a participant,
because it answers "what is collected" and not "what do I do, and how long does
it take". Figure 4 redraws the same content as a grid whose cells state what the
participant does, with a totals row giving the approximate hours each visit
consumes. Both forms are given, because a participant may want to hand the
conventional one to their own oncologist \cite{NCITrial2024}.

\begin{pafloat}
\begin{pafig}[d2-type: grid-rows, grid-columns]
\node[d2cellh,minimum height=9mm,text width=32mm] (h0) at (0,0) {What you do};
\foreach \i/\lab in {1/{Screen\\-28 to -1},2/{Base\\day -1},3/{Surgery\\day 0},4/{Acute\\day 1-7},5/{Day 30},6/{Day 90},7/{q12wk to\\24 mo}}{
  \node[d2cellh,minimum height=9mm,text width=13mm] at ({2.85+(\i-1)*1.62},0) {\lab};}
\foreach \r/\lab in {1/{Sign consent, or decline},2/{Lie still for a scan},3/{Give blood, CA 19-9 and panel},4/{Answer daily-activity questions},5/{Be shown the simulation sign-off},6/{Confirm or withdraw the opt-out},7/{Be asleep under anaesthesia},8/{Have the drain fluid checked},9/{Accept or decline the restart},10/{Report anything that feels wrong},11/{Receive your own results}}{
  \node[d2celll,minimum height=6.4mm,text width=32mm] at (0,{-0.78-(\r-1)*0.66}) {\lab};}
\foreach \r/\c in {1/{2,0,0,0,0,0,0},2/{1,0,0,0,0,1,1},3/{1,1,1,1,1,1,1},4/{1,1,0,0,1,1,1},5/{0,1,0,0,0,0,0},6/{2,2,0,0,0,0,0},7/{0,0,1,0,0,0,0},8/{0,0,0,1,1,0,0},9/{0,0,0,2,0,0,0},10/{1,1,1,1,1,1,1},11/{0,0,0,0,1,1,1}}{
  \foreach \k [count=\j from 1] in \c {
    \ifnum\k=0 \node[d2cell,minimum height=6.4mm,text width=13mm] at ({2.85+(\j-1)*1.62},{-0.78-(\r-1)*0.66}) {};\fi
    \ifnum\k=1 \node[d2cellk,minimum height=6.4mm,text width=13mm] at ({2.85+(\j-1)*1.62},{-0.78-(\r-1)*0.66}) {x};\fi
    \ifnum\k=2 \node[d2cell,fill=protoblue,text=protowhite,minimum height=6.4mm,text width=13mm] at ({2.85+(\j-1)*1.62},{-0.78-(\r-1)*0.66}) {you decide};\fi}}
\node[d2cellg,minimum height=6.4mm,text width=32mm,font=\tiny\sffamily\bfseries] at (0,-8.04) {Hours of your time, approximate};
\foreach \j/\v in {1/{9 h},2/{3 h},3/{6 h},4/{inpatient},5/{2 h},6/{3 h},7/{12 h}}{
  \node[d2cellg,minimum height=6.4mm,text width=13mm] at ({2.85+(\j-1)*1.62},-8.04) {\v};}
\begin{scope}[on background layer]
\node[d2cont,fit={(-1.85,0.62) (12.6,-8.42)}] (G) {};
\end{scope}
\node[d2title,anchor=south west] at (G.north west) {soa: your schedule, activity by visit};
\legkey{-1.7}{-9.1}{protoblue}{a point where you decide}
\legkey{2.3}{-9.1}{pablue1}{done, with you present}
\legkey{6.1}{-9.1}{protowhite}{nothing at this visit}
\legkey{9.7}{-9.1}{pagrayl}{your time cost}
\end{pafig}
\vspace{-0.7cm}
\figcaption{Figure 4. The Schedule of Activities redrawn as a true grid from the participant's\\
chair, with each cell stating what they do rather than what is collected, and a\\
totals row giving the approximate hours of their time each visit will consume.}
\end{pafloat}

\subsection{Schedule of Activities, in the conventional form}

\vspace{0.30em}

\begin{xltabular}{\textwidth}{L{3.9cm} C{1.55cm} C{1.15cm} C{1.05cm} C{1.35cm} C{0.95cm} C{0.95cm} Y}
\toprule
\textbf{Activity} & \textbf{Screen -28 to -1} & \textbf{Base Day -1} & \textbf{Surg Day 0} & \textbf{Acute Day 1-7} & \textbf{Day 30} & \textbf{Day 90} & \textbf{LT q12wk} \\
\midrule
\endfirsthead
\multicolumn{8}{@{}l}{\footnotesize\itshape Table continued from the previous page.}\\
\toprule
\textbf{Activity} & \textbf{Screen -28 to -1} & \textbf{Base Day -1} & \textbf{Surg Day 0} & \textbf{Acute Day 1-7} & \textbf{Day 30} & \textbf{Day 90} & \textbf{LT q12wk} \\
\midrule
\endhead
\midrule
\multicolumn{8}{r@{}}{\footnotesize\itshape Table continued on the next page.}\\
\endfoot
\bottomrule
\endlastfoot
Informed consent, including Physical AI opt-out & X & & & & & & \\
Staging and KRAS G12 confirmation & X & & & & & & \\
CA 19-9 & X & X & & & X & X & X \\
ECOG performance status & X & X & & & X & X & X \\
Safety laboratory panel & X & X & X & X & X & X & X \\
Phase 0 simulation sign-off & & X & & & & & \\
USL at least 7.0 and safety matrix & & X & & & & & \\
Robotic Whipple surgery & & & X & & & & \\
Intra-operative telemetry capture & & & X & & & & \\
ISGPS fistula grading & & & & X & X & & \\
Daraxonrasib restart advisory & & & & X & & & \\
Adverse event review & X & X & X & X & X & X & X \\
RECIST imaging & X & & & & & X & X \\
PFS and OS assessment & & & & & & X & X \\
\end{xltabular}

\vspace{0.30em}

\noindent Abbreviations: ISGPS, International Study Group on Pancreatic Surgery;
LT, long-term; OS, overall survival; PFS, progression-free survival; RECIST,
Response Evaluation Criteria in Solid Tumors; USL, Unified Safety Level.

\subsection{How many people, and when your turn comes}

Up to 18 participants are enrolled across three dose levels, three per level and
expandable to six if a dose-limiting toxicity appears. That number is not a
budget decision. It is what the 3+3 rule needs in order to answer the dose
question, and enrolling more people than the question needs is itself a harm.

Enrollment is deliberately slow, and the reason matters to anyone waiting. A
dose level does not open until the level below it has been cleared, and within a
level the second participant is not dosed until the first participant's sentinel
window has closed. A participant who is screened while a sentinel window is open
waits for it to close. That wait is not an administrative delay and it is not
negotiable; it is the mechanism that makes commitment three true.

The practical consequence is that the calendar in \S 9.1 describes the
participant's own 24 months, not the study's. The study runs longer than any one
participant's participation, and results from participants ahead of you are
reviewed by the Data and Safety Monitoring Board before your level opens. A
participant at DL3 is therefore protected by every safety observation made at
DL1 and DL2, and \S 6.2 states what happens if one of those observations goes
the wrong way while that participant is still enrolled.

There is no waiting list in the ordinary sense, because eligibility is
re-assessed at the moment a slot opens rather than at the moment of the first
call. Disease progresses; a participant who was resectable in March may not be
in June. \S 7.1 states the criteria that are re-checked, and \S 8.1 states what
the site owes a participant who becomes ineligible while waiting, which is a
prompt referral back to their own oncologist with the study's imaging in hand.

\subsection{What the study cannot tell you}

A Phase 1 study of at most eighteen participants is powered for safety and
feasibility and for nothing else. It cannot establish that this combination
improves survival, and no result it produces should be read as though it had.
The National Cancer Institute's guidance on how trials work makes the same point
about early-phase studies generally \cite{NCITrial2024}, and the FDA's public
information on computer-assisted surgical systems states that cancer-related
long-term outcomes have not been established for every use of robotically
assisted surgical devices \cite{FDARobot2022}.

This is stated here, at the front, rather than in a limitations paragraph at the
back, because the surveyed literature is consistent that patients discount
enthusiasm and trust candour \cite{Jauniaux2025, BrarRAS2024X}. A paper arguing
for participation has to be able to say what participation will not deliver.
What this study can deliver is a clean answer to three questions: whether the
dose is tolerated, whether the device causes harm within 30 days, and whether
the operation achieves a clear margin at a rate consistent with the best
contemporary comparator.
